\scanpage{119}
isom qui ne dépend plus de signes. \deleted{Ceci} On peut donc considérer
\[
\left\{\begin{aligned} \delta_H\delta_{H'}&\in(\det B)^{\otimes-2}\simeq D^{\otimes-2}\\ \text{de }\overline{\mathrm m}\qquad\delta_B&\in(\det B)^{\otimes-2}\simeq D^{\otimes-2}
\end{aligned}\right.
\]

Je rappelle que la \deleted{fait} condition que \deleted{\unsure{B}}\deleted{\unsure{B}}$H,H'$ soient ``en position générale'' i.e. engendrent l'\uline{algèbre} $A$, est que
\[
\begin{aligned}
\delta_B\quad&\text{est section \uline{inv} de }(\det\unsure{D|B})^{\otimes-2}\\ &\hspace{7em}\simeq D^{\otimes-2}
\end{aligned}
\]
i.e $D$ anisotrope. Ceci dit, on trouve \margin{si $H,H'$ étales sur $k$ i.e. $\delta_H\delta_{H'}$ section inversible}
\[
\boxed{\rho_0=\frac{\delta_H\delta_{H'}}{\delta_B}}
\]

Cette formule n'est définie pour le moment que si $H,H'$ sont étales /$k$ (et en position générale). On peut la prendre comme \uline{définition} de $\rho_0$ \deleted{\unsure{dès}}\inserted{dès} que $\delta_B$ inv. i.e. $H,H'$ en position générale (pas néc\textsuperscript{t} étale). On peut même définir $\rho_0$ encore plus généralement, du moment que $\delta_H\delta_{H'}$ et $\delta_B$ ne
% archive sheet number: 118
